% Bibliography — 中文翻译（保留原始 \bibitem 格式，每条后附中文标题注释）
% 对应 main.tex 第 2223--2358 行 \begin{thebibliography}...\end{thebibliography}
\begin{thebibliography}{99}
\raggedright
% ============================================================================
\bibitem{ovadia2019calibration} Y.~Ovadia, E.~Fertig, J.~Ren, Z.~Nado,
D.~Sculley, S.~Nowozin, J.~Dillon, B.~Lakshminarayanan, J.~Snoek,
``Can you trust your model's uncertainty? Evaluating predictive
uncertainty under dataset shift,'' \emph{NeurIPS}, 2019.
arXiv:1906.02530.
% 中文标题：你能信任模型的不确定性吗？数据集偏移下的预测不确定性评估
\bibitem{guo2017calibration} C.~Guo, G.~Pleiss, Y.~Sun, K.~Weinberger,
``On calibration of modern neural networks,'' \emph{ICML}, PMLR
70:1321--1330, 2017. arXiv:1706.04599.
% 中文标题：论现代神经网络的校准
\bibitem{wagner2020ptbxl} P.~Wagner, N.~Strodthoff, R.-D.~Bousseljot,
et al., ``PTB-XL, a large publicly available electrocardiography
dataset,'' \emph{Scientific Data}, 7:154, 2020.
doi:10.1038/s41597-020-0495-6.
% 中文标题：PTB-XL，一个大规模公开可获取的心电数据集
\bibitem{zheng2022ecgarrhythmia} J.~Zheng, J.~Guo, H.~Chu,
``A large scale 12-lead electrocardiogram database for arrhythmia
study (Chapman-Shaoxing e-cardiology),'' version 1.0.0,
\emph{PhysioNet}, 2022. doi:10.13026/wgex-er52.
% 中文标题：用于心律失常研究的大规模 12 导联心电图数据库（Chapman-Shaoxing 电子心脏病学）
\bibitem{goldberger2000} A.~L. Goldberger et al., ``PhysioBank,
PhysioToolkit, and PhysioNet: components of a new research resource
for complex physiologic signals,'' \emph{Circulation},
101(23):e215--e220, 2000.
% 中文标题：PhysioBank、PhysioToolkit 与 PhysioNet：面向复杂生理信号新研究资源的组成部分
\bibitem{nixon2019} J.~Nixon, M.~W. Dusenberry, L.~Zhang, et al.,
``Measuring calibration in deep learning,'' \emph{CVPR Workshops},
pp.~38--41, 2019. arXiv:1904.01685.
% 中文标题：深度学习中的校准度量
\bibitem{kull2019dirichlet} M.~Kull, M.~Perello-Nieto, M.~Kangse,
T.~S.~Filho, P.~Song, P.~Flach, ``Beyond temperature scaling:
obtaining well-calibrated multi-class probabilities with Dirichlet
calibration,'' \emph{NeurIPS}, 2019. arXiv:1910.12656.
% 中文标题：超越温度缩放：用 Dirichlet 校准获取良好校准的多类概率
\bibitem{alexandari2020} A.~Alexandari, A.~Kundaje, A.~Shrikumar,
``Maximum likelihood with bias-corrected calibration is hard-to-beat
at label shift adaptation,'' \emph{ICML}, PMLR 119:222--232, 2020.
arXiv:1901.06852.
% 中文标题：带偏差校正校准的最大似然在标签偏移适应中难以被超越
\bibitem{lipton2018bbs} Z.~Lipton, Y.-X.~Wang, A.~Smola,
``Detecting and correcting for label shift with black box
predictors,'' \emph{ICML}, PMLR 80:1944--1953, 2018.
arXiv:1802.03916.
% 中文标题：用黑盒预测器检测和校正标签偏移
\bibitem{barandas2023} M.~Barandas, L.~Famiglini, A.~Campagner,
D.~Folgado, R.~Sim\~ao, F.~Cabitza, H.~Gamboa, ``Evaluation of
uncertainty quantification methods in multi-label classification: a
case study with automatic diagnosis of electrocardiogram,''
\emph{Information Fusion}, 101:101978, 2024.
doi:10.1016/j.inffus.2023.101978.
% 中文标题：多标签分类中不确定性量化方法的评估：以心电图自动诊断为案例研究
\bibitem{zhang2022} J.~F.~Vranken et al., ``Uncertainty estimation for
deep learning-based automated analysis of 12-lead
electrocardiograms,'' \emph{European Heart Journal -- Digital Health},
2(3):423--434, 2021. doi:10.1093/ehjdh/ztab045.
% 中文标题：基于深度学习的 12 导联心电图自动分析的不确定性估计
\bibitem{reyna2022} M.~Reyna, N.~Strodthoff, P.~Wagner, et al.,
``Issues in the automated classification of multilead ECGs using
heterogeneous labels and populations,'' \emph{Physiological
Measurement}, 43(8):084001, 2022. doi:10.1088/1361-6579/ac79fd.
% 中文标题：使用异构标签和人群进行多导联心电图自动分类中的若干问题
\bibitem{efron1987bca} B.~Efron, ``Better bootstrap confidence
intervals,'' \emph{Journal of the American Statistical Association},
82(397):171--185, 1987. doi:10.1080/01621459.1987.10478410.
% 中文标题：更好的自助法置信区间
\bibitem{vancalster2016} B.~Van Calster, D.~J.~Steyerberg,
S.~Collins, R.~D.~Riley, E.~W.~Steyer, ``A calibration hierarchy for
risk models was defined: from utopia to empirical data,''
\emph{Journal of Clinical Epidemiology}, 74:167--176, 2016.
doi:10.1016/j.jclinepi.2015.12.005.
% 中文标题：风险模型校准层级的定义：从乌托邦到经验数据
\bibitem{vancalster2019} B.~Van Calster, D.~J.~McLernon,
M.~van~Smeden, L.~Wynants, G.~S.~Collins, ``Calibration: the Achilles
heel of predictive analytics,'' \emph{BMC Medicine}, 17:230, 2019.
doi:10.1186/s12916-019-1466-7.
% 中文标题：校准：预测分析的阿喀琉斯之踵
\bibitem{kumar2019} A.~Kumar, P.~Liang, T.~Ma, ``Verified uncertainty
calibration,'' \emph{NeurIPS}, 2019. arXiv:1909.10155.
% 中文标题：经验证的不确定性校准
\bibitem{strodthoff2021} N.~Strodthoff, P.~Wagner, T.~Schaeffter,
W.~Samek, ``Deep learning for ECG analysis: benchmarks and insights
from PTB-XL,'' \emph{IEEE Journal of Biomedical and Health
Informatics}, 25(5):1519--1528, 2021.
doi:10.1109/JBHI.2020.3022989.
% 中文标题：用于心电图分析的深度学习：来自 PTB-XL 的基准与见解
\bibitem{physiolmeas2026} M.~Haekal, S.~N.~Khotimah,
G.~R.~F.~Suwandi, F.~Haryanto, ``RR-interval-based atrial fibrillation
detection and burden estimation: cross-dataset validation and
calibration-aware probability analysis,'' \emph{Physiological
Measurement}, 2026. doi:10.1088/1361-6579/ae99aa.
% 中文标题：基于 RR 间期的心房颤动检测与负荷估计：跨数据集验证与校准感知概率分析
\bibitem{medrxiv2026} K.~Patel, P.~Beedala, ``Calibration drift under
cross-institutional deployment: an external validation framework for
ICU mortality prediction across MIMIC-IV and eICU,'' \emph{medRxiv}
(preprint), 2026. doi:10.64898/2026.05.03.26352335.
% 中文标题：跨机构部署下的校准漂移：跨 MIMIC-IV 与 eICU 的 ICU 死亡率预测外部验证框架
\bibitem{transcal2020} X.~Wang, M.~Long, J.~Wang, M.~I.~Jordan,
``TransCal: transferability estimation for auto-selecting transfer
calibration,'' \emph{NeurIPS}, 2020. arXiv:2007.08259.
% 中文标题：TransCal：用于自动选择迁移校准的可迁移性估计
\bibitem{pseudocal2020} S.~Garg, Y.~Wu, S.~Balakrishnan, Z.~Lipton,
``A unified view of label shift estimation,'' \emph{NeurIPS}, 2020.
arXiv:2003.07554.
% 中文标题：标签偏移估计的统一视角
\bibitem{lascal2024} T.~Popordanoska, G.~Radevski, T.~Tuytelaars, et al.,
``LaSCal: label-shift calibration for uncertainty quantification,''
\emph{NeurIPS}, 2024.
% 中文标题：LaSCal：用于不确定性量化的标签偏移校准
\bibitem{morenotorres2012} J.~T.~Moreno-Torres, T.~Raeder,
R.~Alaiz-Rodriguez, N.~Chawla, F.~Herrera, ``A unifying view on
dataset shift in classification,'' \emph{Pattern Recognition},
45(1):521--530, 2012. doi:10.1016/j.patcog.2011.05.019.
% 中文标题：分类中数据集偏移的统一视角
\bibitem{protocolv21a1} [Unpublished pre-registered protocol] ``ECG
calibration boundary study -- pre-registered experiment protocol
v2.1-A1,'' 2026. Hash manifest:
\texttt{docs/osf\_archive\_manifest.json} (local snapshot; public OSF
archival pending); revision A1 full text in
\texttt{docs/PROTOCOL\_AMENDMENT\_A1\_PREREG\_RELOCATION.md}.
% 中文标题：心电图校准边界研究——预注册实验协议 v2.1-A1
\bibitem{nosek2019prereg} B.~Nosek, C.~Ebersole, A.~DeHaven,
D.~Mellor, ``The preregistration revolution,'' \emph{PNAS},
115(11):2600--2606, 2018. doi:10.1073/pnas.1708274114.
% 中文标题：预注册革命
\bibitem{lakens2019prereg} D.~Lakens, ``The practical alternative to
the $p$ value is the correctly used $p$ value,''
\emph{Perspectives on Psychological Science}, 16(3):639--648, 2021.
doi:10.1177/1745691620958012.
% 中文标题：$p$ 值的实用替代品是正确使用的 $p$ 值
\bibitem{ecgfounder2024} J.~Li, A.~Aguirre, V.~M.~Junior, et al.,
``ECGFounder: a general-purpose foundation model for electrocardiogram
analysis trained on over 10 million recordings,'' \emph{NEJM AI},
1(6):AIoa2401033, 2024. arXiv:2410.04133.
% 中文标题：ECGFounder：在超过 1000 万条记录上训练的心电图分析通用基础模型
\bibitem{ecgfm2024} J.~C.~McHugh, B.~Wang, et al., ``ECG-FM: an open
electrocardiogram foundation model,'' \emph{arXiv preprint}
2408.05178, 2024.
% 中文标题：ECG-FM：一个开放的心电图基础模型
\bibitem{heartlang2025} J.~Jin, et al., ``Reading your heart: learning
ECG words and sentences via pre-training ECG language model
(HeartLang),'' \emph{ICLR}, 2025. arXiv:2502.10707.
% 中文标题：读懂你的心：通过预训练心电图语言模型（HeartLang）学习心电词与句子
\bibitem{transferecg2025} C.~V.~Nguyen, C.~D.~Do, ``Transfer learning in ECG
diagnosis: is it effective?'' \emph{PLOS ONE}, 20(5):e0316043, 2025.
doi:10.1371/journal.pone.0316043.
% 中文标题：心电图诊断中的迁移学习：它有效吗？
\bibitem{adaptood2026} S.~Vavaroutas, et al., ``ADAPTOOD:
uncertainty-aware fine-tuning for out-of-distribution ECG time series
models,'' \emph{arXiv preprint} 2606.04164, 2026.
% 中文标题：ADAPTOOD：面向分布外心电图时间序列模型的感知不确定性的微调
\bibitem{beatrhythm2026} W.~Jiang, et al., ``BeatRhythm-TTA: test-time
adaptation for ECG classification via SQI-gated self-training and
beat-rhythm consistency,'' \emph{arXiv preprint} 2608.23347, 2026.
% 中文标题：BeatRhythm-TTA：通过 SQI 门控自训练与心搏节律一致性进行心电图分类的测试时适应
\bibitem{star2025} N.~Nemati, ``Comparative analysis of data
augmentation for clinical ECG classification with STAR (Sinusoidal
Time--Amplitude Resampling),'' \emph{arXiv preprint} 2510.24740, 2025.
% 中文标题：使用 STAR（正弦时间-幅度重采样）进行临床心电图分类数据增强的比较分析
\bibitem{peace2026} X.~Liu, et al., ``PEACE: knowledge-guided
cross-modal fusion for adult-to-pediatric ECG transfer via
label-conditioned contrastive alignment,'' \emph{arXiv preprint}
2605.00647, 2026.
% 中文标题：PEACE：通过标签条件对比对齐实现成人到儿童心电图迁移的知识引导跨模态融合
\bibitem{song2024ecgfm} J.~Song, J.-H.~Jang, B.~T.~Lee, D.~Hong,
J.~Kwon, Y.-Y.~Jo, ``Foundation models for electrocardiograms,''
\emph{arXiv preprint} 2407.07110, 2024.
% 中文标题：心电图基础模型
\end{thebibliography}
